% Part: lambda-calculus
% Chapter: church-rosser
% Section: beta-reduction

\documentclass[../../../include/open-logic-section]{subfiles}

\begin{document}

\olfileid[id]{lam}{cr}{b}

\olsection{Reduksi-$\beta$}

\begin{lem}\ollabel{lem:one-par}
  Jika $M \bredone M'$, maka $M \bredpar M'$.
\end{lem}
\begin{proof} Jika $M \bredone M'$, maka untuk suatu $x$, $N$, dan~$Q$,
    $M=(\lambd[x][N])Q$ dan $M'=\Subst{N}{Q}{x}$. Karena
    $N\bredpar N$ dan $Q\bredpar Q$ berdasarkan \olref[pb]{thm:refl},
    kita langsung memperoleh $(\lambd[x][N])Q
    \bredpar\Subst{N}{Q}{x}$ berdasarkan
    \olref[pb]{defn:bredpar}\olref[pb]{defn:bredpar4}.
\end{proof}

\begin{lem}\ollabel{lem:par-red}
  Jika $M \bredpar M'$, maka $M \bred M'$.
\end{lem}

\begin{proof} Dengan induksi pada !!{derivation} $M \bredpar M'$.
  \begin{enumerate}
  \item Aturan terakhir adalah \olref[pb]{defn:bredpar1}: Maka $M$ dan $M'$
    keduanya adalah $x$, dan $x \bred x$.
  \item Aturan terakhir adalah \olref[pb]{defn:bredpar2}: Untuk suatu $x$,
    $N$, dan~$N'$, $M=\lambd[x][N]$, $M'=\lambd[x][N']$, dan
    $N\bredpar N'$. Berdasarkan hipotesis induksi, $N\bred N'$. Maka
    $\lambd[x][N]\bred\lambd[x][N']$ melalui rangkaian kontraksi
    $\bredone$ yang sama seperti dalam $N\bred N'$.
  \item Aturan terakhir adalah \olref[pb]{defn:bredpar3}: Untuk suatu $P$,
    $Q$, $P'$, dan~$Q'$, $M=PQ$, $M'=P'Q'$, $P\bredpar P'$, dan
    $Q\bredpar Q'$. Berdasarkan hipotesis induksi, $P\bred P'$ dan
    $Q\bred Q'$. Jadi $PQ\bred P'Q'$ dengan terlebih dahulu mereduksi
    $P\bred P'$, kemudian $Q\bred Q'$.
  \item Aturan terakhir adalah \olref[pb]{defn:bredpar4}: Untuk suatu $x$,
    $N$, $N'$, $Q$, dan~$Q'$, $M=(\lambd[x][N])Q$,
    $M'=\Subst{N'}{Q'}{x}$, $N\bredpar N'$, dan $Q\bredpar Q'$.
    Berdasarkan hipotesis induksi, $Q\bred Q'$ dan $N\bred N'$. Jadi
    $(\lambd[x][N])Q\bred\Subst{N'}{Q'}{x}$ dengan mereduksi $N\bred N'$,
    kemudian $Q\bred Q'$, dan akhirnya mengontraksi
    $(\lambd[x][N'])Q'$ menjadi $\Subst{N'}{Q'}{x}$.
  \end{enumerate}
\end{proof}

\begin{lem}\ollabel{lem:str}
  $\bred$ adalah relasi transitif terkecil yang memuat $\bredpar$.
\end{lem}

\begin{proof}
  Misalkan $\xred$ adalah relasi transitif terkecil yang memuat $\bredpar$.

  $\bred \subseteq \xred$: Andaikan $M \bred M'$, yaitu
  $M\ident M_1\bredone\dots\bredone M_k\ident M'$. Berdasarkan
  \olref{lem:one-par}, $M\ident M_1\bredpar\dots\bredpar M_k\ident M'$.
  Karena $\xred$ memuat $\bredpar$ dan transitif, $M\xred M'$.

  $\xred \subseteq \bred$: Andaikan $M \xred M'$, yaitu
  $M\ident M_1\bredpar\dots\bredpar M_k\ident M'$. Berdasarkan
  \olref{lem:par-red}, $M\ident M_1\bred\dots\bred M_k\ident M'$. Karena
  $\bred$ transitif, $M\bred M'$.
\end{proof}

\begin{thm}\ollabel{thm:cr}
  $\bred$ memenuhi sifat Church--Rosser.
\end{thm}

\begin{proof}
  Langsung dari \olref[dap]{thm:str}, \olref[pb]{thm:cr}, dan
  \olref{lem:str}.
\end{proof}

\end{document}
